
\clearpage
\pageoriginale

\SGAExposeBookmark{XI}{Cohomología de las intersecciones completas}
\begin{center}
{\large\textbf{SGA 7}}\\[1.5em]
{\Large\underline{EXPOSÉ XI}}\\[2.5em]
{\Large\underline{COHOMOLOGÍA DE LAS INTERSECCIONES COMPLETAS}}\\[1.5em]
{\large\underline{por P. DELIGNE}}
\end{center}

\bigskip
\noindent\underline{1. Cohomología del espacio proyectivo y de las
intersecciones completas}

Sea $X$ una variedad proyectiva compleja no singular, puramente de dimensión
$n$, y sea $\underline O(1)$ un haz invertible amplio sobre $X$. Según
Akizuki y Nakano [1], del teorema de anulación de Kodaira [4] se deduce que

\begin{equation}\tag{*}
H^i(X,\Omega_X^j(-1))=0\qquad\text{para }i+j<n.
\end{equation}

También según [1], este resultado implica el «teorema de Lefschetz débil»
en cohomología compleja. Nos proponemos demostrar, por fuerza bruta, que si
$X$ es un espacio proyectivo sobre un cuerpo de característica arbitraria,
entonces (*) sigue siendo cierto. Esto permitirá, a imitación de [1],
comparar las cohomologías de Hodge y de De Rham de las intersecciones
completas, sobre un cuerpo de característica arbitraria, con las del espacio
proyectivo.

A continuación explicitaremos algunas consecuencias de los resultados
obtenidos. En el n.º 2 exponemos resultados cualitativos deducidos de la
fórmula numérica de Hirzebruch (2.2). Quiero agradecer a N. Katz la ayuda que
me prestó en la concepción de ese número.

\medskip
\noindent 1.0. Sea $\underline E$ un módulo localmente libre de rango $r+1$
sobre un esquema $S$. \underline{Suponemos que $r>0$.} Recordemos que
$\mathbb P(\underline E)$ designa el espacio proyectivo relativo sobre $S$
de los hiperplanos de $\underline E$. Designaremos por $p$ la proyección de
$\mathbb P(\underline E)$ sobre $S$, y por

\begin{equation*}
\eta\in H^0\bigl(S,R^1p_*\Omega^1_{\mathbb P(\underline E)/S}\bigr)
\end{equation*}

la primera clase de Chern, en el sentido de Hodge, del haz invertible
$\underline O(1)$.

\pageoriginale

Recordemos que $\eta$ es la imagen, por la diferencial logarítmica

\begin{equation*}
df/f:\underline O^*\longrightarrow
\Omega^1_{\mathbb P(\underline E)/S},
\end{equation*}

de la clase de $\underline O(1)$ en $H^0(S,R^1p_*\underline O^*)$.
Para todo haz $\underline F$ de $\underline O$-módulos sobre
$\mathbb P(\underline E)$, ponemos

\begin{equation*}
\underline F(n)=\underline F\otimes\underline O(1)^{\otimes n}
\qquad(n\in\mathbb Z).
\end{equation*}

\medskip
\noindent\underline{Teorema 1.1.}\quad
(i) \underline{Los haces
$R^ip_*(\Omega^j_{\mathbb P(\underline E)/S}(n))$ son localmente libres, y
su formación es compatible con todo cambio de base.}

(ii) \underline{Para $0\leq i\leq r$, la sección $\eta^i$ de
$R^ip_*(\Omega^i_{\mathbb P(\underline E)/S})$ define un isomorfismo entre
este haz y $\underline O_S$. Además,}

\begin{equation*}
R^ip_*\Omega^j_{\mathbb P(\underline E)/S}=0
\qquad\underline{\text{si }}i\ne j\ \underline{\text{ o }}\ i>r.
\end{equation*}

(iii) \underline{Si $n\ne0$, entonces, fuera de los dos casos siguientes,
$R^ip_*(\Omega^j_{\mathbb P(\underline E)/S}(n))$ es nulo:}

\begin{align}
i&=0\quad\underline{\text{y}}\quad n>j,
  \tag{iii$_a$}\\
i&=r\quad\underline{\text{y}}\quad n<j-r.
  \tag{iii$_b$}
\end{align}

Sean
\[
\mathbb V(\underline E)=
\operatorname{Spec}\bigl(\operatorname{Sym}^*\underline E\bigr)
\]
el espacio afín definido por $\underline E$, $e$ su sección nula,
$\mathbb V(\underline E)^*=\mathbb V(\underline E)\setminus e(S)$, $a$ y
$a'$ las proyecciones de $\mathbb V(\underline E)$ y
$\mathbb V(\underline E)^*$ sobre $S$, y $q$ la proyección de
$\mathbb V(\underline E)^*$ sobre $\mathbb P(\underline E)$:

\begin{center}
\begin{tikzcd}[row sep=large,column sep=large]
\mathbb V(\underline E)
  \arrow[dr,"a"]
& \mathbb V(\underline E)^*
  \arrow[l,hook']
  \arrow[r,"q"]
  \arrow[d,"a'"']
& \mathbb P(\underline E)
  \arrow[dl,"p"]
\\
{} & S\arrow[ul,bend right,"e"'] & {}
\end{tikzcd}
\qquad .
\end{center}

El espacio afín punteado $\mathbb V(\underline O(1))^*$ sobre
$\mathbb P(\underline E)$ es canónicamente isomorfo a
$\mathbb V(\underline E)^*$, de modo que, para todo haz cuasicoherente
$\underline F$ de $\underline O$-módulos sobre $\mathbb P(\underline E)$,
se tiene canónicamente

\pageoriginale

\begin{equation}\tag{1.1.1}
q_*q^*\underline F\simeq
\bigoplus_{n\in\mathbb Z}\underline F(n).
\end{equation}

Esta graduación de $q_*q^*\underline F$ puede deducirse de la acción, sobre
este haz, del grupo $\mathbb G_m$ de las homotecias de
$\mathbb V(\underline E)$.

Los campos de vectores relativos sobre $\mathbb V(\underline E)$ se
identifican con los $S$-morfismos de $\mathbb V(\underline E)$ en
$\mathbb V(\underline E)$, identificado con el espacio tangente en el
origen. El campo de vectores $X$ definido por la aplicación identidad de
$\mathbb V(\underline E)$ es invariante por homotecias. Este campo es
tangente a las fibras de $q$ y define una trivialización

\begin{equation}\tag{1.1.2}
X:\Omega^1_{\mathbb V(\underline E)^*/\mathbb P(\underline E)}
\longrightarrow\underline O_{\mathbb V(\underline E)^*}.
\end{equation}

Si los $x_i$ forman una base de $\underline E$, se tiene

\begin{equation*}
X=\sum_i x_i\,dx^i.
\end{equation*}

Como $\Omega^1_{\mathbb V(\underline E)^*/\mathbb P(\underline E)}$ es de
rango uno, la sucesión exacta

\begin{equation*}
0\longrightarrow q^*\Omega^1_{\mathbb P(\underline E)/S}
\longrightarrow\Omega^1_{\mathbb V(\underline E)^*/S}
\longrightarrow
\Omega^1_{\mathbb V(\underline E)^*/\mathbb P(\underline E)}
\longrightarrow0
\end{equation*}

induce sucesiones exactas

\begin{equation*}
0\longrightarrow q^*\Omega^i_{\mathbb P(\underline E)/S}
\longrightarrow\Omega^i_{\mathbb V(\underline E)^*/S}
\longrightarrow
\Omega^1_{\mathbb V(\underline E)^*/\mathbb P(\underline E)}
\otimes q^*\Omega^{i-1}_{\mathbb P(\underline E)/S}
\longrightarrow0.
\end{equation*}

Mediante (1.1.2), estas se reescriben

\begin{equation}\tag{1.1.3}
0\longrightarrow q^*\Omega^i_{\mathbb P(\underline E)/S}
\longrightarrow\Omega^i_{\mathbb V(\underline E)^*/S}
\xrightarrow{\ \iota_X\ }
q^*\Omega^{i-1}_{\mathbb P(\underline E)/S}
\longrightarrow0.
\end{equation}

Canónicamente,

\begin{equation}\tag{1.1.4}
\Omega^i_{\mathbb V(\underline E)/S}
=a^*\bigwedge^i\underline E.
\end{equation}

Las sucesiones exactas cortas (1.1.3) se obtienen, por tanto, fraccionando
la siguiente sucesión exacta larga sobre $\mathbb V(\underline E)^*$:

\begin{equation}\tag{1.1.5}
\cdots\longrightarrow {a'}^*\bigwedge^i\underline E
\xrightarrow{\ \iota_X\ }
{a'}^*\bigwedge^{i-1}\underline E
\xrightarrow{\ \iota_X\ }
{a'}^*\bigwedge^{i-2}\underline E\longrightarrow\cdots.
\end{equation}

\pageoriginale

Esta sucesión exacta larga no es otra cosa que la restricción a
$\mathbb V(\underline E)^*$ de la resolución de Koszul de
$\underline O_{e(S)}$ sobre $\mathbb V(\underline E)$.

Proyectemos la sucesión exacta larga (1.1.5) sobre
$\mathbb P(\underline E)$ mediante el morfismo afín $q$. En virtud de
(1.1.1) y de la fórmula $q_*{a'}^*=q_*q^*p^*$, se obtiene una sucesión
exacta larga

\begin{equation}\tag{1.1.6}
\cdots\longrightarrow
\bigoplus_{n\in\mathbb Z}p^*\bigwedge^i\underline E(n)
\xrightarrow{\ \iota_X\ }
\bigoplus_{n\in\mathbb Z}p^*\bigwedge^{i-1}\underline E(n)
\longrightarrow\cdots .
\end{equation}

Esta se fracciona en sucesiones exactas cortas, imágenes de (1.1.3),

\begin{equation}\tag{1.1.7}
0\longrightarrow
\bigoplus_{n\in\mathbb Z}\Omega^i_{\mathbb P(\underline E)/S}(n)
\longrightarrow
\bigoplus_{n\in\mathbb Z}p^*\bigwedge^i\underline E(n)
\xrightarrow{\ \iota_X\ }
\bigoplus_{n\in\mathbb Z}\Omega^{i-1}_{\mathbb P(\underline E)/S}(n)
\longrightarrow0.
\end{equation}

Las sucesiones exactas (1.1.7) son sumas directas de las sucesiones exactas

\begin{equation}\tag{1.1.8}
0\longrightarrow\Omega^i_{\mathbb P(\underline E)/S}(n)
\longrightarrow p^*\bigwedge^i\underline E(n-i)
\xrightarrow{\ \iota_X\ }
\Omega^{i-1}_{\mathbb P(\underline E)/S}(n)
\longrightarrow0.
\end{equation}

Recordemos (EGA III, 2.1.15, o FAC) que los haces
$R^ip_*\underline O(k)$ son localmente libres y su formación es compatible
con todo cambio de base, y que

\begin{equation}\tag{1.1.9}
R^ip_*\underline O(k)=0
\quad\text{si }i\ne0,r,
\quad\text{o si }i=0,\ k<0,
\quad\text{o si }i=r,\ k>-r-1,
\end{equation}

y que, por consiguiente, $R^ip_*p^*\bigwedge^j\underline E(k)$ es nulo bajo
las mismas condiciones.

Las sucesiones exactas (1.1.8) proporcionan una resolución

\begin{equation}\tag{1.1.10}
0\longrightarrow\Omega^i_{\mathbb P(\underline E)/S}(n)
\xrightarrow{\ E\ }p^*\bigwedge^i\underline E(n-i)
\xrightarrow{\ \iota_X\ }\cdots\longrightarrow
p^*\bigwedge^0\underline E(n)\longrightarrow0
\end{equation}

de $\Omega^i_{\mathbb P(\underline E)/S}(n)$. La sucesión espectral
definida por esta resolución degenera ($E_1=E_\infty$) por razones de grado,
según (1.1.9); para $0<j<r$, proporciona

\begin{equation}\tag{1.1.11}
R^jp_*\Omega^i_{\mathbb P(\underline E)/S}(n)
=\mathcal H\left(\cdots\xrightarrow{\ \iota_X\ }
p_*\bigwedge^{i-j}\underline E(n-i+j)
\xrightarrow{\ \iota_X\ }\cdots\right).
\end{equation}

Como $r\geq1$, un argumento de profundidad muestra que

\begin{equation}\tag{1.1.12}
\bigoplus_{n\in\mathbb Z}p_*p^*\bigwedge^i\underline E(n)
=p_*q_*q^*p^*\bigwedge^i\underline E
={a'}_*{a'}^*\bigwedge^i\underline E
=a_*a^*\bigwedge^i\underline E.
\end{equation}

\pageoriginale

Por tanto, $p_*p^*\bigwedge^i\underline E(n)$ es la parte homogénea de peso
$n+i$ de $a_*a^*\bigwedge^i\underline E$, es decir,
$\operatorname{Sym}^n(\underline E)\otimes\bigwedge^i\underline E$.
Recordando que (1.1.5) procede de la resolución de Koszul, se deduce de
(1.1.11) que, para $0<j<r$, el haz
$R^jp_*\Omega^i_{\mathbb P(\underline E)/S}(n)$ solo puede ser no nulo si
$i-j=0$ y $n-i+j=0$, es decir, si $n=0$ e $i=j$. En este caso, además,
$R^ip_*\Omega^i_{\mathbb P(\underline E)/S}$ es isomorfo a
$\underline O_S$.

Para $j=0$, se deduce también de (1.1.10) que

\begin{equation}\tag{1.1.13}
p_*\Omega^i_{\mathbb P(\underline E)/S}(n)
=\ker\!\left(\iota_X:
p_*\bigwedge^i\underline E(n-i)\longrightarrow
p_*\bigwedge^{i-1}\underline E(n-i+1)\right),
\end{equation}

y el mismo argumento muestra que $\iota_X$ es inyectivo cuando $n-i\leq0$.

En virtud de (1.1.8), para $n=0$ e $i=r+1$, se tiene

\begin{equation*}
\Omega^r_{\mathbb P(\underline E)/S}\simeq
p^*\bigwedge^{r+1}\underline E(-r-1).
\end{equation*}

Cuando $S$ es el espectro de un cuerpo, las afirmaciones de anulación de
1.1(ii) y (iii) resultan de lo anterior para $j\ne r$, y se deducen por
dualidad ([2], V, 11.2, p.~211) para $j=r$. En particular, cualesquiera que
sean $j$ y $n$, existe a lo sumo un $i$ tal que
$R^ip_*\Omega^j_{\mathbb P(\underline E)/S}(n)\ne0$. Esto implica (i),
puesto que $\Omega^j_{\mathbb P(\underline E)/S}$ es plano sobre $S$ (lo
más difícil es encontrar la referencia en EGA III; sugiero: 6.10.5,
7.4.1 y 7.5.5, seguidos de un paso al límite).

Sabiendo que los $R^ip_*\Omega^i_{\mathbb P(\underline E)/S}$ son haces
invertibles para $0\leq i\leq r$, solo queda probar que, cuando $S$ es el
espectro de un cuerpo, $\eta^i\ne0$. En efecto, $\eta^r$ es la clase de
cohomología de un punto racional (intersección de $r$ hiperplanos), y por
tanto una base de $H^r(\mathbb P^r,\Omega^r)$ (FGA, Exposé 149, n.º 4).

\medskip
\noindent 1.2. Es posible calcular directamente a partir de (1.1.10) o
(1.1.13) las características de Euler--Poincaré de los $\Omega^i(n)$. El
resultado, poco apetitoso, es

\begin{equation*}
\chi\bigl(\mathbb P^r_k,\Omega^j_{\mathbb P^r_k}(n)\bigr)
=\sum_{0\leq i\leq j}(-1)^{i-j}
\binom{r+1}{i}\binom{r+n-i}{r}.
\end{equation*}

\pageoriginale

\medskip
\noindent\underline{Proposición 1.3.}\quad
\underline{Sea $f:X\to S$ un morfismo propio y liso, sea $\underline W$ un
haz localmente libre de rango $c$ sobre $X$, sea $\underline F$ un haz
coherente sobre $X$, sea $d>0$ un entero, y sea
$s:\underline W\to\underline O_X$ una sección de $\underline W^\vee$.
Supongamos que $S$ es noetheriano y que:}

\begin{enumerate}[label=(\alph*),leftmargin=2.4em]
\item \underline{El subesquema cerrado $H$ de $X$ definido por la anulación
de $s$ es liso sobre $S$.}

\item \underline{Localmente sobre $X$, si los $s_i$ son las coordenadas de
$s$ en una base de $\underline W^\vee$, los $s_i$ forman a la vez una
sucesión $\underline O_X$-regular y una sucesión
$\underline F$-regular.}

\item \underline{Para toda sucesión de enteros $k_\ell$, no todos nulos, se
tiene}
\end{enumerate}

\begin{equation*}
R^if_*\left(\bigotimes_\ell\bigwedge^{k_\ell}\underline W
\otimes\Omega^j_{X/S}\otimes\underline F\right)=0
\qquad\underline{\text{para }}i+j<d.
\end{equation*}

\underline{Entonces:}

\begin{align*}
R^if_*\left(\Omega^j_{\sisi{H}{X}/S}\otimes\underline F\right)
&\xrightarrow{\ \sim\ }
R^if_*\left(\Omega^j_{\sisi{X}{H}/S}\otimes\underline F\right)
&&\underline{\text{para }}i+j<d-c, \tag{i}\\
R^if_*\left(\Omega^j_{\sisi{H}{X}/S}\otimes\underline F\right)
&\lhook\joinrel\longrightarrow
R^if_*\left(\Omega^j_{\sisi{X}{H}/S}\otimes\underline F\right)
&&\underline{\text{para }}i+j=d-c. \tag{ii}
\end{align*}

El haz conormal $N_{H/X}$ de $H$ en $X$ es isomorfo a la restricción de
$\underline W$ a $H$, de modo que la sucesión exacta

\begin{equation}\tag{1.3.1}
0\longrightarrow N_{H/X}\longrightarrow
\Omega^1_{X/S}\otimes\underline O_H\longrightarrow
\Omega^1_{H/S}\longrightarrow0
\end{equation}

define una filtración decreciente sobre
$\Omega^j_{X/S}\otimes\underline O_H$ cuyos cocientes sucesivos son

\begin{equation}\tag{1.3.2}
\operatorname{Gr}^k\bigl(\Omega^j_{X/S}\otimes\underline O_H\bigr)
\simeq\bigwedge^k\underline W\otimes\Omega^{j-k}_{H/S}.
\end{equation}

De (b) se sigue que también se tiene

\begin{equation}\tag{1.3.3}
\operatorname{Gr}^k\bigl(\Omega^j_{X/S}\otimes\underline O_H
\otimes\underline F\bigr)
=\bigwedge^k\underline W\otimes\Omega^{j-k}_{H/S}
\otimes\underline F.
\end{equation}

Demostramos (i) y (ii) por inducción sobre $j$, simultáneamente para todos
los $\underline F$ que satisfacen (b) y (c). Las afirmaciones son vacías
para $j<0$.

\pageoriginale

Supongamos, pues, que (i) y (ii) se han demostrado para $j'<j$. La
hipótesis de inducción se aplica a los haces
$\bigwedge^k\underline W\otimes\underline F$ y, por (c), da

\begin{equation}\tag{1.3.4}
R^if_*\left(\bigwedge^k\underline W\otimes
\Omega^{j-k}_{H/S}\otimes\underline F\right)=0
\quad\text{para }k>0\text{ e }i+j<d-c+k.
\end{equation}

La filtración de
$\Omega^j_{X/S}\otimes\underline O_H\otimes\underline F$ considerada
arriba define una sucesión espectral

\begin{equation*}
E_1^{p,i-p}=R^if_*\left(\bigwedge^p\underline W\otimes
\Omega^{j-p}_{H/S}\otimes\underline F\right)
\Longrightarrow
R^if_*\left(\Omega^j_{X/S}\otimes\underline O_H
\otimes\underline F\right).
\end{equation*}

Por (1.3.4), $E_1^{p,i-p}=0$ cuando $p>0$ e $i+j\leq d-c$; por tanto,

\begin{equation}\tag{1.3.5}
\left\{
\begin{aligned}
R^if_*\left(\Omega^j_{X/S}\otimes\underline O_H\otimes\underline F\right)
&\xrightarrow{\ \sim\ }
R^if_*\left(\Omega^j_{H/S}\otimes\underline F\right)
&&\text{si }i+j<d-c,\\
R^if_*\left(\Omega^j_{X/S}\otimes\underline O_H\otimes\underline F\right)
&\lhook\joinrel\longrightarrow
R^if_*\left(\Omega^j_{H/S}\otimes\underline F\right)
&&\text{si }i+j=d-c.
\end{aligned}
\right.
\end{equation}

Por la hipótesis (b), con una diferencial apropiada, los haces
$\bigwedge^k\underline W\otimes\Omega^j_{X/S}\otimes\underline F$ forman
una resolución de Koszul por la izquierda de
$\Omega^j_{X/S}\otimes\underline O_H\otimes\underline F$. Esta filtración
define una sucesión espectral

\begin{equation}\tag{1.3.6}
E_1^{-p,q}=R^qf_*\left(\bigwedge^p\underline W\otimes
\Omega^j_{X/S}\otimes\underline F\right)
\Longrightarrow
R^{q-p}f_*\left(\Omega^j_{X/S}\otimes\underline O_H
\otimes\underline F\right).
\end{equation}

Por la hipótesis (c), $E_1^{-p,q}=0$ para $p>0$ y $q+j<d$, y trivialmente
para $p\notin[0,c]$. En consecuencia,

\begin{equation}\tag{1.3.7}
\left\{
\begin{aligned}
R^if_*\left(\Omega^j_{X/S}\otimes\underline F\right)
&\xrightarrow{\ \sim\ }
R^if_*\left(\Omega^j_{H/S}\otimes\underline F\right)
&&\text{si }i+j<d-c,\\
R^if_*\left(\Omega^j_{X/S}\otimes\underline F\right)
&\lhook\joinrel\longrightarrow
R^if_*\left(\Omega^j_{H/S}\otimes\underline F\right)
&&\text{si }i+j=d-c.
\end{aligned}
\right.
\end{equation}

Las afirmaciones (i) y (ii) se deducen ahora de (1.3.5) y (1.3.7).

\pageoriginale

\medskip
\noindent 1.4. Sea $\underline E$ un haz localmente libre de rango $r+1$
sobre un esquema $S$, sea $\mathbb P(\underline E)$ el fibrado proyectivo
correspondiente, sea $d\leq r$ un entero y sea
$\underline a=(a_1,\ldots,a_d)$ una sucesión de enteros. Una
\underline{intersección completa} relativa de multigrado $\underline a$
en $\mathbb P(\underline E)$ es un subesquema cerrado que tiene codimensión
$d$ en cada fibra y que, localmente sobre $S$ para la topología étale, es la
intersección de $d$ hipersuperficies de grados respectivos
$a_1,\ldots,a_d$.

Por 1.1, las hipótesis (b) y (c) de 1.3 se satisfacen para
$X/S=\mathbb P(\underline E)/S$, $\underline F=\underline O_X$,

\begin{equation*}
\underline W=\bigoplus_i\underline O(-a_i),
\end{equation*}

y una sección $s$ de $\underline W^\vee$ que define una intersección
completa relativa $Y\subset\mathbb P(\underline E)$ de multigrado
$\underline a$.

\medskip
\noindent\underline{Teorema 1.5.}\quad
\underline{Bajo las hipótesis de 1.4, sea $n=r-d$ la dimensión relativa de
$Y$, y sea $\eta\in H^0(S,R^1f_*\Omega^1_{Y/S})$ la primera clase de Chern
de $\underline O(1)$ en el diagrama}

\begin{center}
\begin{tikzcd}[column sep=large,row sep=large]
Y\arrow[rr,hook]\arrow[dr,dash,"f"']&&
\mathbb P(\underline E)\arrow[dl,"g"]\\
&S&{}
\end{tikzcd}
\qquad .
\end{center}

\underline{Si $Y$ es liso sobre $S$, entonces:}

\begin{enumerate}[label=(\roman*),leftmargin=2.6em]
\item \underline{Los haces $R^if_*\Omega^j_{Y/S}$ son localmente libres,
y su formación es compatible con todo cambio de base.}

\item \underline{La sucesión espectral}
\begin{equation*}
E_1^{p,q}=R^qf_*\Omega^p_{Y/S}\Longrightarrow
R^{p+q}f_*(\Omega^*_{Y/S}),
\end{equation*}
\underline{que relaciona las cohomologías de Hodge y de De Rham, degenera
($E_1=E_\infty$); su abutimiento es localmente libre y su formación es
compatible con todo cambio de base.}
\end{enumerate}

\noindent(iii)$_a$ \underline{Se tiene}

\begin{equation*}
R^if_*\Omega^j_{Y/S}=0
\qquad\underline{\text{si }}i\ne j\ \underline{\text{ e }}\ i+j\ne n.
\end{equation*}

\pageoriginale

\noindent(iii)$_b$ \underline{$\eta^i$ es una base de
$R^if_*\Omega^i_{Y/S}$ si $2i<n$.}

\noindent(iii)$_c$ \underline{Si $n<2i\leq2n$, y si $\alpha_i$ es la base
de $R^if_*\Omega^i_{Y/S}$ dual de la base $\eta^{n-i}$ de
$R^{n-i}f_*\Omega^{n-i}_{Y/S}$, entonces}

\begin{equation*}
\eta^i=\left(\prod_{h=1}^d a_h\right)\alpha_i.
\end{equation*}

\noindent(iii)$_d$ \underline{Si $n=2m$, entonces $\eta^m\ne0$ (siempre
que $S\ne\varnothing$), y el cociente
$R^mf_*\Omega^m_{Y/S}/\underline O_S\eta^m$ es localmente libre.}

Basta demostrar 1.5 en un caso «universal», lo que permite suponer que $S$
es liso sobre $\operatorname{Spec}(\mathbb Z)$. En este caso, (ii) se
deduce de (i), pues vale sobre $\mathbb C$, y por tanto sobre todo cuerpo de
característica cero, por la teoría de Hodge. Además, como $S$ es reducido,
(i) se deduce de (iii), supuesto demostrado cuando $S$ es el espectro de un
cuerpo. En efecto, por (iii), para cada $j$ las funciones

\begin{equation*}
s\longmapsto\dim_{k(s)}H^i(Y_s,\Omega^j_{Y_s})
\end{equation*}

son constantes sobre $S$, salvo quizá cuando $i=j$; por consiguiente, todas
son constantes, ya que su suma alternada es constante (EGA III, 7.9.2). La
conclusión se deduce entonces de EGA III, 6.10.5, 7.6.9 y 7.7.5.

Por 1.4 y 1.3, se tiene

\begin{equation}\tag{1.5.1}
\left\{
\begin{aligned}
R^ig_*\Omega^j_{\mathbb P(\underline E)/S}
&\xrightarrow{\ \sim\ }R^if_*\Omega^j_{Y/S}
&&\text{para }i+j<n,\\
R^ig_*\Omega^j_{\mathbb P(\underline E)/S}
&\lhook\joinrel\longrightarrow R^if_*\Omega^j_{Y/S}
&&\text{para }i+j=n.
\end{aligned}
\right.
\end{equation}

Cuando $S$ es el espectro de un cuerpo, las afirmaciones (iii)$_a$,
(iii)$_b$ y (iii)$_d$ se deducen de (1.5.1) y de la dualidad de Serre
([2], V, 11.2, p.~211). La clase de cohomología

\pageoriginale

$\eta^n$ es la clase de un $0$-ciclo de grado
$\prod_{h=1}^d a_h$ (el grado de la intersección completa), de modo que

\begin{equation*}
\eta^n=\left(\prod_{h=1}^d a_h\right)\alpha_n.
\end{equation*}

Por definición, $\eta^{n-i}\alpha_i=\alpha_n$, y de ello se deduce
(iii)$_c$.

Así quedan demostrados (i) y (ii); en el caso universal, las afirmaciones de
(iii) se deducen de (i) y de (iii) en el caso ya demostrado, en el que $S$
es el espectro de un cuerpo.

\medskip
\noindent\underline{Teorema 1.6.}\quad
\underline{Conservemos las hipótesis y notaciones de 1.4 y 1.5, supongamos
que $Y$ es liso, sea $\ell$ un primo invertible en todos los cuerpos
residuales de $S$, y sea
$\eta_\ell\in H^0(S,R^2f_*\mathbb Z_\ell(1))$ la primera clase de Chern
$\ell$-ádica de $\underline O(1)$ (SGA 5, VII). Entonces:}

\begin{enumerate}[label=(\roman*),leftmargin=2.6em]
\item \underline{$R^if_*\mathbb Z_\ell=0$ si $i$ es impar e $i\ne n$.}

\item \underline{Si $2i<n$, entonces $R^{2i}f_*\mathbb Z_\ell(i)$ es
isomorfo a $\mathbb Z_\ell$ y tiene a $\eta_\ell^i$ como base.}

\item \underline{Si $n<2i\leq2n$, entonces
$R^{2i}f_*\mathbb Z_\ell(i)$ es isomorfo a $\mathbb Z_\ell$ y tiene}
\begin{equation*}
\eta_\ell^i\bigg/\prod_{h=1}^d a_h
\end{equation*}
\underline{como base.}

\item \underline{El $\mathbb Z_\ell$-haz $R^nf_*\mathbb Z_\ell$ es
constructible, torcido constante y sin torsión. Si $n$ es par, entonces
$\eta_\ell^{n/2}\ne0$ (siempre que $S\ne\varnothing$), y
$R^nf_*\mathbb Z_\ell(n/2)/\mathbb Z_\ell\eta_\ell^{n/2}$ es sin torsión.}
\end{enumerate}

Los $\mathbb Z_\ell$-haces $R^if_*\mathbb Z_\ell$ son constructibles y
torcidos constantes, y su formación es compatible con todo cambio de base,
por SGA 4, XVI, 2.2. Por tanto, basta demostrar 1.6 cuando $S$ es el
espectro de un cuerpo algebraicamente cerrado. La aplicación repetida de
SGA 4, XIV, 3.3 muestra entonces, para todo $m$, que la aplicación canónica

\begin{equation*}
H_Y^q\bigl(\mathbb P(\underline E),\mathbb Z/\ell^m\bigr)
\longrightarrow H^q(Y,\mathbb Z/\ell^m)
\end{equation*}

\pageoriginale

es biyectiva para $q>r+d$ y sobreyectiva para $q=r+d$. De esto y de la
dualidad de Poincaré (SGA 4, XVIII), se deduce (SGA 5, VII) que la aplicación
de restricción

\begin{equation*}
H^q\bigl(\mathbb P(\underline E),\mathbb Z_\ell\bigr)
\longrightarrow H^q(Y,\mathbb Z_\ell)
\end{equation*}

es biyectiva para $q<n$, e inyectiva con conúcleo sin torsión para $q=n$.
Esto demuestra (i) para $i<n$, así como (ii) y (iv). Como en 1.5, se
comprueba que $\eta_\ell^n$ es $\prod_{h=1}^d a_h$ veces la clase canónica;
(i) y (iii) se deducen entonces de la dualidad de Poincaré (SGA 4, XVIII).

\medskip
\noindent 1.7. Si $Y$ es una intersección completa de dimensión al menos
tres en un espacio proyectivo $\mathbb P^r(k)$ sobre un cuerpo $k$, se sabe
que $\operatorname{Pic}(Y)$ es isomorfo a $\mathbb Z$ y está engendrado por
$\underline O(1)$ (SGA 2, XII, 3.7). En dimensión dos, cuando $Y$ es lisa,
todavía puede demostrarse lo siguiente.

\medskip
\noindent\underline{Teorema 1.8.}\quad
\underline{Sea $Y$ una superficie lisa, intersección completa en el espacio
proyectivo $\mathbb P^r(k)$ sobre un cuerpo algebraicamente cerrado $k$.
Entonces $\underline{\operatorname{Pic}}^\tau(Y)=0$, y la clase de
$\underline O(1)$ no es divisible en el grupo de Néron--Severi
$\operatorname{NS}(Y)$.}

Como $H^1(Y,\underline O)=0$ por 1.5(iii)$_a$, se sabe que
$\underline{\operatorname{Pic}}^0(Y)=0$. La anulación de
$\underline{\operatorname{Pic}}^\tau(Y)$ se deduce entonces del lema
familiar siguiente, del hecho de que $H^2(Y,\mathbb Z_\ell(1))$ es sin
torsión por 1.6(iv), y del hecho de que $H^0(Y,\Omega_Y^1)=0$ por
1.5(iii)$_a$.

\medskip
\noindent\underline{Lema 1.9.}\quad
\underline{Sea $Y$ una variedad algebraica completa sobre un cuerpo
algebraicamente cerrado $k$.}

\begin{enumerate}[label=(\roman*),leftmargin=2.6em]
\item \underline{Si $\ell$ es primo con la característica de $k$, entonces
la $\ell$-torsión de $\operatorname{NS}(Y)$ se identifica con la
$\ell$-torsión de $H^2(Y,\mathbb Z_\ell(1))$.}

\pageoriginale

\item \underline{Si $k$ tiene característica no nula $p$ e $Y$ es lisa,
entonces el núcleo de la multiplicación por $p$ en $\operatorname{Pic}(Y)$
se identifica con el subgrupo de $H^0(Y,\Omega_Y^1)$ formado por las
diferenciales localmente logarítmicas.}
\end{enumerate}

Las sucesiones exactas de Kummer

\begin{equation*}
0\longrightarrow\mu_{\ell^n}\longrightarrow\mathbb G_m
\xrightarrow{\ x\mapsto x^{\ell^n}\ }\mathbb G_m\longrightarrow0
\end{equation*}

dan sucesiones exactas

\begin{equation*}
0\longrightarrow
\operatorname{Pic}(Y)/\ell^n\operatorname{Pic}(Y)
\longrightarrow H^2(Y,\mu_{\ell^n})
\longrightarrow H^2(Y,\mathbb G_m)_{\ell^n}
\longrightarrow0,
\end{equation*}

donde

\begin{equation*}
\operatorname{Pic}(Y)/\ell^n\operatorname{Pic}(Y)
=\operatorname{NS}(Y)\otimes\mathbb Z/\ell^n.
\end{equation*}

Pasando al límite inverso se obtiene

\begin{equation*}
0\longrightarrow\operatorname{NS}(Y)\otimes\mathbb Z_\ell
\longrightarrow H^2(Y,\mathbb Z_\ell(1))
\longrightarrow
\operatorname{Hom}\bigl(\mathbb Q_\ell/\mathbb Z_\ell,
H^2(Y,\mathbb G_m)\bigr)
\longrightarrow0.
\end{equation*}

El conúcleo es sin torsión porque $\mathbb Q_\ell/\mathbb Z_\ell$ es
divisible. Esto demuestra (i).

Para demostrar (ii), basta observar que la sucesión exacta

\begin{equation*}
0\longrightarrow\underline O_Y^*
\xrightarrow{\ x\mapsto x^p\ }\underline O_Y^*
\xrightarrow{\ df/f\ }\Omega_Y^1
\end{equation*}

define una sucesión exacta corta

\begin{equation*}
0\longrightarrow\underline O_Y^*\longrightarrow\underline O_Y^*
\longrightarrow\Omega^1_{Y,\mathrm{loc}\,\log}\longrightarrow0,
\end{equation*}

y tomar después la sucesión exacta larga de cohomología correspondiente.
Este argumento se debe a Cartier, quien demuestra además que, si $C$
designa el operador de Cartier sobre $H^0(Y,\Omega_Y^1)$, las diferenciales
localmente logarítmicas forman el núcleo de $C-\operatorname{Id}$.

Supongamos, para llegar a una contradicción, que la clase de
$\underline O(1)$ en $\operatorname{NS}(Y)$ es de la forma $nu$, con
$n>1$. Si $\ell$ divide a $n$ y es primo con la característica, entonces en
cohomología $\ell$-ádica se tiene de nuevo

\pageoriginale

\begin{equation*}
\eta_\ell=n\,c_1(u),
\end{equation*}

lo que contradice 1.6(iv).

Finalmente, si la característica $p$ de $k$ es no nula y $p$ divide a $n$,
entonces en cohomología de Hodge se tiene de nuevo

\begin{equation*}
\eta=n\,c_1(u)=0,
\end{equation*}

lo que contradice 1.5(iii)$_d$ y termina la demostración.

\bigskip
\noindent\underline{2. Resultados numéricos}

\medskip
\noindent 2.1. Sea $\underline a=(a_i)_{1\leq i\leq d}$ una sucesión de
enteros $\geq1$, y sea $V_n(\underline a)$ una intersección completa lisa de
dimensión $n$ y multigrado $\underline a$ en un espacio proyectivo
$\mathbb P^{n+d}$ sobre un cuerpo $k$. Pongamos

\begin{equation*}
h^{p,q}(\underline a)
=\dim_k H^q\bigl(V_{p+q}(\underline a),\Omega^p\bigr)
\end{equation*}

y

\begin{equation*}
h_0^{p,q}(\underline a)=h^{p,q}(\underline a)-\delta_{p,q}.
\end{equation*}

Así, los $h_0^{p,q}(\underline a)$ son los números de Hodge de la
cohomología primitiva de dimensión media de una intersección completa de
multigrado $\underline a$. Dependen únicamente de $\underline a$, $p$ y
$q$, por 1.5. Pongamos

\begin{equation}\tag{2.1.1}
H(\underline a)=\sum_{p,q\geq0}h_0^{p,q}(\underline a)y^pz^q
\in\mathbb Z[[y,z]].
\end{equation}

Para $d=1$ y $\underline a=(a)$, pongamos
$h^{p,q}(a)=h^{p,q}(\underline a)$ y $H(a)=H(\underline a)$. Para $d=0$,
pongamos $\sup(a_i)=1$.

La fórmula siguiente es el caso particular $k=0$ de la fórmula (2),
p.~160, de Hirzebruch [3]:

\pageoriginale

\begin{equation}\tag{2.2}
\sum_{n=0}^{\infty}\sum_p
\chi\bigl(V_n(\underline a),\Omega^p\bigr)y^pz^{n+d}
=\frac{1}{(1+yz)(1-z)}
\prod_{i=1}^d
\frac{(1+yz)^{a_i}-(1-z)^{a_i}}
{(1+yz)^{a_i}z+y(1-z)^{a_i}}.
\end{equation}

La transformaremos en la fórmula siguiente.

\medskip
\noindent\underline{Teorema 2.3 (Hirzebruch).}\quad
\underline{Con las notaciones de 2.1, se tiene}

\begin{equation*}
H(\underline a)=\frac{1}{(1+y)(1+z)}\left[
\prod_{i=1}^d
\frac{(1+y)^{a_i}-(1+z)^{a_i}}
{-(1+y)^{a_i}z+(1+z)^{a_i}y}-1\right].
\end{equation*}

Por 1.5,

\begin{equation*}
\chi\bigl(V_n(\underline a),\Omega^p\bigr)
=(-1)^p+(-1)^{n-p}h_0^{p,n-p}(\underline a).
\end{equation*}

Por consiguiente, el miembro izquierdo $I$ de (2.2) puede reescribirse

\begin{align*}
I
&=\sum_{n=0}^{\infty}\sum_{p+q=n}
\left((-1)^p+(-1)^q h_0^{p,q}(\underline a)\right)y^pz^{p+q+d}\\
&=\sum_{p,q\geq0}(-1)^p(yz)^pz^qz^d
+\sum_{p,q\geq0}(-1)^q h_0^{p,q}(\underline a)(yz)^pz^qz^d\\
&=z^d\left[\frac{1}{(1+yz)(1-z)}
+H(\underline a)(yz,-z)\right].
\end{align*}

El miembro derecho $II$ de (2.2) puede reescribirse

\begin{equation*}
II=\frac{z^d}{(1+yz)(1-z)}
\prod_{i=1}^d
\frac{(1+yz)^{a_i}-(1-z)^{a_i}}
{(1+yz)^{a_i}z+yz(1-z)^{a_i}}.
\end{equation*}

Tras simplificar $z^d$ en $I=II$, se obtiene

\begin{equation*}
H(\underline a)(yz,-z)
=\frac{1}{(1+yz)(1-z)}\left[
\prod_{i=1}^d
\frac{(1+yz)^{a_i}-(1-z)^{a_i}}
{(1+yz)^{a_i}z+(1-z)^{a_i}yz}-1\right],
\end{equation*}

y 2.3 se deduce mediante un cambio de variables.

\pageoriginale

\medskip
\noindent\underline{Corolario 2.4.}\quad
\underline{Con las notaciones de 2.1:}

\begin{enumerate}[label=(\roman*),leftmargin=2.6em]
\item \underline{Para $d=1$, se tiene}
\begin{equation*}
\begin{aligned}
H(a)
&=-\frac{(1+y)^{a-1}-(1+z)^{a-1}}
{(1+y)^az-(1+z)^ay}\\[1ex]
&=\frac{\displaystyle\sum_{i,j\geq0}
\binom{a-1}{i+j+1}y^iz^j}
{\displaystyle1-\sum_{i,j\geq1}\binom{a}{i+j}y^iz^j}.
\end{aligned}
\end{equation*}

\item \underline{Si $|P|$ designa el número de elementos de un subconjunto
$P$ de $[1,d]$, se tiene}
\begin{equation*}
H(\underline a)=
\sum_{\substack{P\subset[1,d]\\P\ne\varnothing}}
(1+y)^{|P|-1}(1+z)^{|P|-1}
\prod_{i\in P}H(a_i).
\end{equation*}
\end{enumerate}

Por 2.3,

\begin{align*}
H(a)
&=-\frac{1}{(1+y)(1+z)}\left[
\frac{(1+y)^a-(1+z)^a}{(1+y)^az-(1+z)^ay}+1\right]\\
&=-\frac{1}{(1+y)(1+z)}\left[
\frac{(1+y)^a(1+z)-(1+z)^a(1+y)}
{(1+y)^az-(1+z)^ay}\right]\\
&=-\frac{(1+y)^{a-1}-(1+z)^{a-1}}
{(1+y)^az-(1+z)^ay}.
\end{align*}

Sean $N$ y $D$ el numerador y el denominador de esta última expresión, de
modo que $H(a)=-N/D$.

\pageoriginale

Se tiene

\begin{equation*}
N=\sum_k\binom{a-1}{k}(y^k-z^k)
=(y-z)\sum_{i,j\geq0}\binom{a-1}{i+j+1}y^iz^j
\end{equation*}

y

\begin{align*}
D
&=\sum_k\binom ak(y^kz-yz^k)
=(z-y)+yz\sum_{k\geq1}\binom ak(y^{k-1}-z^{k-1})\\
&=(y-z)\left(-1+yz\sum_{i,j\geq0}
\binom a{i+j+2}y^iz^j\right)\\
&=(y-z)\left(-1+\sum_{i,j\geq1}
\binom a{i+j}y^iz^j\right).
\end{align*}

Esto demuestra (i).

Invirtamos ahora la fórmula de 2.3 cuando $d=1$. Se obtiene

\begin{equation*}
-\frac{(1+y)^a-(1+z)^a}{(1+y)^az-(1+z)^ay}
=1+(1+y)(1+z)H(a).
\end{equation*}

Sustituyendo esta identidad en 2.3 se obtiene

\begin{align*}
H(\underline a)
&=\frac{1}{(1+y)(1+z)}\left[
\prod_{i=1}^d\left(1+(1+y)(1+z)H(a_i)\right)-1\right]\\
&=\frac{1}{(1+y)(1+z)}\left[
\sum_{\substack{P\subset[1,d]\\P\ne\varnothing}}
(1+y)^{|P|}(1+z)^{|P|}\prod_{i\in P}H(a_i)\right],
\end{align*}

lo que demuestra (ii).

\medskip
\noindent\underline{Teorema 2.5.}\quad
\underline{Con las notaciones de 2.1:}

\begin{enumerate}[label=(\roman*),leftmargin=2.6em]
\item \underline{Si $p+q=p'+q'$ y $p\leq p'\leq q'\leq q$, entonces}
\begin{equation*}
h_0^{p,q}(\underline a)\leq h_0^{p',q'}(\underline a).
\end{equation*}

\item \underline{Si $p\leq q$, entonces $h_0^{p,q}(\underline a)\ne0$ si
y solo si}
\begin{equation*}
p>\left\lfloor
\frac{p+q+d-\sum_i a_i}{\sup_i(a_i)}
\right\rfloor.
\end{equation*}
\underline{Aquí $\lfloor\ \rfloor$ designa la parte entera del cociente.}
\end{enumerate}

\pageoriginale

Sea (*) la propiedad siguiente de una serie formal
$P=\sum a_{p,q}y^pz^q$:

\begin{quote}
(*) los $a_{p,q}$ son no negativos, $a_{p,q}=a_{q,p}$, y siempre que
$p+q=p'+q'$ y $p\leq p'\leq q'\leq q$, se tiene
$a_{p,q}\leq a_{p',q'}$.
\end{quote}

La afirmación (i) se deduce inmediatamente de las fórmulas de 2.4, de la
identidad $(1-x)^{-1}=\sum_{n\geq0}x^n$ y del lema siguiente, cuya
demostración se deja al lector.

\medskip
\noindent\underline{Lema 2.5.1.}\quad
\underline{El conjunto de las series formales que satisfacen (*) es
estable por suma y multiplicación.}

Demostremos primero (ii) en el caso particular de las hipersuperficies
($d=1$). Para $p\leq q$, la fórmula buscada puede reescribirse

\begin{equation}\tag{2.5.2}
h_0^{p,q}(a)\ne0
\quad\Longleftrightarrow\quad
p\geq\left\lfloor\frac{p+q+1}{a}\right\rfloor
\end{equation}

o, equivalentemente, $(p+1)a\geq p+q+2$. De la segunda fórmula de 2.4(i)
se ve que $h_0^{p,q}(a)$ es no nulo si y solo si existen pares $(e,f)$ y
$(c_i,d_i)$ tales que

\begin{equation}\tag{2.5.3}
\left\{
\begin{aligned}
(e,f)&\geq(0,0),& e+f&\leq a-2,\\
(c_i,d_i)&\geq(1,1),& c_i+d_i&\leq a,\\
(p,q)&=(e,f)+\sum_i(c_i,d_i).&&
\end{aligned}
\right.
\end{equation}

\pageoriginale

Para $p+q$ fijo, el menor $p$ posible se obtiene tomando $e=0$ y todos los
$c_i=1$; es el número de índices $i$. Es, por tanto, la cota inferior de
los $p$ para los cuales

\begin{equation*}
p\geq\frac{(p+q)-(a-2)}{a},
\qquad\text{es decir,}\qquad (p+1)a\geq p+q+2.
\end{equation*}

Pasemos ahora al caso general. Si $\underline a'$ se obtiene de
$\underline a$ suprimiendo los $a_i$ iguales a uno, entonces
$H(\underline a)=H(\underline a')$. En vista de la forma de 2.5(ii),
podemos suponer, por tanto, que todos los $a_i>1$. Si $d=0$, entonces
$h_0^{p,q}=0$ y la afirmación es correcta; no se pierde generalidad al
suponer $d>0$ y $a_d=\sup_i(a_i)$.

Por 2.4(ii) y (2.5.2), la condición
$h_0^{p,q}(\underline a)\ne0$ equivale a la existencia de un subconjunto no
vacío $P$ de $[1,d]$, un par $(e,f)$ y pares
$(c_i,d_i)_{i\in P}$ tales que

\begin{equation}\tag{2.5.4}
\left\{
\begin{aligned}
(|P|-1,|P|-1)&\geq(e,f)\geq(0,0),\\
c_i,d_i&\geq
\left\lfloor\frac{c_i+d_i+1}{a_i}\right\rfloor,
& (c_i,d_i)&\geq(0,0),\\
(p,q)&=(e,f)+\sum_{i\in P}(c_i,d_i).&&
\end{aligned}
\right.
\end{equation}

Para $n=p+q$ fijo, busquemos el menor valor posible de $p$. Pongamos
$n_i=c_i+d_i$. Para tal $p$, las condiciones (2.5.4) se convierten en

\pageoriginale

\begin{equation}\tag{2.5.5}
\left\{
\begin{aligned}
a_i=2&\Longrightarrow n_i\text{ es par}
&&\text{(de lo contrario (2.5.4) es imposible)},\\
c_i&\geq\left\lfloor\frac{n_i+1}{a_i}\right\rfloor,
&c_i&\leq n_i,\\
(|P|-1,|P|-1)&\geq(e,f)\geq(0,0),&&\\
n&=e+f+\sum_{i\in P}n_i,&&\\
p&=e+\sum_{i\in P}c_i.&&
\end{aligned}
\right.
\end{equation}

Esta descripción muestra primero que el menor $p$ se obtiene para
$P=[1,d]$, caso al que nos restringimos en adelante. Si

\begin{equation}\tag{2.5.6}
\sum_i(a_i-2)+(d-1)\geq n,
\end{equation}

existe una solución con $p=c_i=e=0$. En caso contrario, si $a_d\ne2$, la
mejor elección es $n_i=a_i-2$ para $i\ne d$, de modo que $c_i=0$ para
$i\ne d$, $e=0$, $f=d-1$, y

\begin{equation*}
n_d=n-(d-1)-\sum_{i\ne d}(a_i-2).
\end{equation*}

Entonces

\begin{equation}\tag{2.5.7}
p=\left\lfloor
\frac{n-(d-1)-\sum_{i\ne d}(a_i-2)+1}{a_d}
\right\rfloor,
\end{equation}

y, por tanto, para $p\leq q$, se tiene
$h_0^{p,q}(\underline a)\ne0$ si y solo si

\begin{equation}\tag{2.5.8}
p\geq\left\lfloor
\frac{n-\sum_{i\ne d}(a_i-1)+1}{a_d}
\right\rfloor.
\end{equation}

La fórmula (2.5.8) se reduce a (2.5.2) para $d=1$ y es trivialmente cierta
cuando se cumple (2.5.6). Si $d\geq2$, (2.5.6) no se cumple y todos los
$a_i$ son iguales a dos, la mejor elección es en cambio

\pageoriginale

$n_i=0$ para $i\ne d$, $e=0$, y $f=d-1$ o $d-2$, elegido congruente con
$n$ módulo dos, con $n_d=n-f$. Si $n\equiv d-1\pmod2$, esto demuestra
(2.5.7) como antes. En caso contrario, obsérvese que

\begin{equation*}
\left\lfloor\frac{n-(d-2)+1}{2}\right\rfloor
=\left\lfloor\frac{n-(d-1)+1}{2}\right\rfloor,
\end{equation*}

de modo que (2.5.7), y por tanto (2.5.8), sigue siendo válida cuando
$a_d=2$.

Queda verificar que (2.5.8) equivale a 2.5(ii). En efecto,

\begin{equation*}
\left\lfloor\frac{n-\sum_{i\ne d}(a_i-1)+1}{a_d}\right\rfloor
=\left\lfloor\frac{n+d-\sum_{i\ne d}a_i}{a_d}\right\rfloor
=\left\lfloor\frac{p+q+d-\sum_i a_i}{a_d}\right\rfloor+1.
\end{equation*}

\medskip
\noindent 2.6. Si $X$ es una intersección completa de dimensión $n$ en
$\mathbb P^r$, se deduce inmediatamente de (2.5.8) que la parte primitiva
de $H^n(X)$ solo se anula cuando $X$ es una cuádrica de dimensión impar.

\medskip
\noindent\underline{Definición 2.7.}\quad
\underline{Si $X$ es una variedad algebraica propia y lisa sobre
$\mathbb C$, el nivel de Hodge de $H^n(X)$ es el supremo de los números
$q-p$ para los cuales $p+q=n$ y $h^{p,q}(X)\ne0$.}

Por tanto, el nivel de Hodge $\ell$ es congruente con $n$ módulo dos y es a
lo sumo
\[
\min\bigl(n,\sisi{\dim(X)-n}{2\dim(X)-n}\bigr).
\]
Es no negativo salvo si $H^n(X)=0$, en cuyo caso vale $-\infty$.

\medskip
\noindent\underline{Corolario 2.8.}\quad
\underline{Conservemos las notaciones de 2.1, pongamos $k=\mathbb C$ y
supongamos que $a_d=\sup_i(a_i)$. Sea $\ell$ un entero tal que
$0\leq\ell<n$, congruente con $n$ módulo dos. El nivel de Hodge de
$H^n(V_n(\underline a))$ es a lo sumo $\ell$ si y solo si}

\pageoriginale

\begin{equation*}
\sum_{i\ne d}(a_i-1)
\leq(\ell+1)-(a_d-2)\frac{n-\ell}{2}.
\end{equation*}

Por (2.5.8), si el nivel de Hodge es no negativo, es el máximo de cero y

\begin{equation*}
n-2\left\lfloor
\frac{n-\sum_{i\ne d}(a_i-1)+1}{a_d}
\right\rfloor.
\end{equation*}

La desigualdad

\begin{equation*}
n-2\left\lfloor
\frac{n-\sum_{i\ne d}(a_i-1)+1}{a_d}
\right\rfloor\leq\ell
\end{equation*}

puede reescribirse

\begin{equation*}
\frac{n-\ell}{2}
\leq
\left\lfloor
\frac{n-\sum_{i\ne d}(a_i-1)+1}{a_d}
\right\rfloor,
\end{equation*}

o, equivalentemente,

\begin{equation*}
\sum_{i\ne d}(a_i-1)
\leq1+n-a_d\frac{n-\ell}{2}
=\ell+1-(a_d-2)\frac{n-\ell}{2}.
\end{equation*}

\medskip
\noindent 2.9. Mediante la fórmula 2.8 es fácil, para $\ell$ pequeño,
enumerar los multigrados de las intersecciones completas en
$\mathbb P^r$ de dimensión $n>1$ cuyo $H^n$ tiene nivel $\ell$. Sin
embargo, la tabla que figuraba en la primera versión de esta exposición era
incorrecta; la tabla siguiente se extrae de la más completa dada por
M. Rapoport en su artículo «Complément à l'article de P. Deligne
\emph{La conjecture de Weil pour les surfaces K3}», \emph{Inventiones
Math.} \underline{15}, 227--236 (1972). Para $\ell\leq1$, las
intersecciones completas $V_n(\underline a)$ de dimensión $n$ y multigrado
$\underline a=(a_1,\ldots,a_d)$, con $n>\ell$, $d\geq1$ y $a_i\geq2$, que
tienen nivel de Hodge $\ell$, son las siguientes:

\pageoriginale

\begin{equation*}
\begin{array}{ll}
\ell=-\infty & V_{2k+1}(2) \qquad (k\geq0),\\[2ex]
\ell=0 & V_{2k}(2)\ (k\geq1),\quad
          V_{2k}(2,2)\ (k\geq1),\quad V_2(3),\\[2ex]
\ell=1 & V_{2k+1}(2,2)\ (k\geq1),\quad
          V_{2k+1}(2,2,2)\ (k\geq1),\\
& V_3(3),\quad V_5(3),\quad V_3(3,2),\quad V_3(4).
\end{array}
\end{equation*}

\medskip
\noindent 2.10. Para $\underline a$ pequeño, las funciones generatrices de
los $h_0^{p,q}(\underline a)$ son, por 2.4,

\begin{equation*}
\begin{array}{rcl}
H(2)&=&\dfrac{1}{1-yz},\\[3ex]
H(3)&=&\dfrac{2+y+z}{1-3yz-y^2z-yz^2},\\[3ex]
H(2,2)&=&\dfrac{y+z}{(1-yz)^2}
          +\dfrac{2}{(1-yz)^2}+\dfrac{1}{1-yz},
\end{array}
\end{equation*}

de modo que

\begin{equation*}
\begin{array}{rcll}
h_0^{p,p}(2)&=&1&,\\[2ex]
h_0^{1,1}(3)&=&6&,\\[2ex]
h_0^{p,p}(2,2)&=&2p+3&.
\end{array}
\end{equation*}

\pageoriginale

\bigskip
\begin{center}
\underline{BIBLIOGRAFÍA}
\end{center}

\begin{enumerate}[label={[\arabic*]},leftmargin=3em]
\item Y. Akizuki y S. Nakano, «Note on Kodaira--Spencer's proof of
Lefschetz theorems», \emph{Proc. Jap. Acad.} 30, 266--272 (1954).

\item R. Hartshorne, \emph{Residues and Duality}, Lecture Notes in
Mathematics 20, Springer, 1966.

\item F. Hirzebruch, \emph{Topological Methods in Algebraic Geometry},
3.ª ed., Springer-Verlag, 1966.

\item K. Kodaira, «On a differential geometric method in the theory of
analytic stacks», \emph{Proc. Nat. Acad. Sci. USA} 39, 1268--1273
(1953).
\end{enumerate}
